% GENERATED FILE. Edit canonical lessons, then run npm run generate:books.
% canonical-source-hash: fnv1a64:a13f65bc6c2799b5
% canonical-lessons: HI-C431-bhav, HI-C431-mol, HI-C431-sauda, HI-C431-sel, HI-C431-kifayati

\chapter{Rates and Deals}
\label{ch:hi-rates-and-deals}
\hlchaptermodality{431}

\begin{chapteropening}
\textbf{By the end of this chapter:} \emph{I can ask the going rate, talk about worth, strike a deal and find a sale or good value.}

Complete the last lesson of chapter 431: I can ask the going rate, talk about worth, strike a deal and find a sale or good value.
\end{chapteropening}

\section[\texorpdfstring{\hi{भाव} (bhāv)}{भाव (bhāv)}]{\hi{भाव} (bhāv) — the going rate; a price per unit}

\label{lesson:HI-C431-bhav}



\begin{quote}
Before the new one: say the Hindi for to be paid off, then the Hindi for to collect (dues).
\end{quote}

\subsection*{You'll want to know: \hi{भाव}}
\textbf{\hi{भाव}} — \emph{bhāv} — \textquotedblleft{}the going rate; a price per unit\textquotedblright{}.

\begin{quote}
\hi{आज} \hi{आलू} \hi{का} \hi{भाव} \hi{क्या} \hi{है}?
\end{quote}

\begin{quote}
\emph{āj ālū kā bhāv kyā hai?}
\end{quote}

\begin{quote}
What is the price of potatoes today?
\end{quote}

Masculine.

\begin{cousinweb}
\textbf{\hi{भाव}} (\emph{bhāv}) is Sanskrit \emph{bhāva}, \textquotedblleft{}being, a state, a condition\textquotedblright{} (Monier-Williams); in Hindi markets it came to mean the state of prices on the day, a sense the dictionary does not give. The same word also means a feeling.
\end{cousinweb}

\begin{grammarlens}[title={What you've built}]
One more word for getting around.
\end{grammarlens}

\subsection*{Your turn}
\begin{itemize}
\raggedright
  \item \emph{Say it:} \emph{bhāv}
  \item \emph{Say it:} \emph{bhāv}, once more
  \item \emph{Say it:} \emph{āj ālū kā bhāv kyā hai?} — What is the price of potatoes today?
  \item \emph{Recall:} say the Hindi for to be paid off, then the Hindi for to collect (dues), then say \emph{bhāv} again
\end{itemize}

\begin{tcolorbox}[breakable,colback=teal!4,colframe=teal!35!black,title={Before you move on}]
What does \hi{भाव} mean? (\textquotedblleft{}The going rate; a price per unit\textquotedblright{}.) Say it once more.
\end{tcolorbox}
\section[\texorpdfstring{\hi{मोल} (mol)}{मोल (mol)}]{\hi{मोल} (mol) — worth; a price (old-fashioned)}

\label{lesson:HI-C431-mol}



\begin{quote}
Before the new one: say the Hindi for to collect (dues), then the Hindi for the going rate.
\end{quote}

\subsection*{You'll want to know: \hi{मोल}}
\textbf{\hi{मोल}} — \emph{mol} — \textquotedblleft{}worth; a price (old-fashioned)\textquotedblright{}.

\begin{quote}
\hi{माँ} \hi{के} \hi{प्यार} \hi{का} \hi{कोई} \hi{मोल} \hi{नहीं} \hi{होता।}
\end{quote}

\begin{quote}
\emph{māṁ ke pyār kā koī mol nahīṃ hotā.}
\end{quote}

\begin{quote}
A mother's love is beyond price.
\end{quote}

Masculine.

\begin{cousinweb}
\textbf{\hi{मोल}} (\emph{mol}) is a tadbhava of Sanskrit \emph{maulya}, \textquotedblleft{}a price\textquotedblright{} (CDIAL 10373), which lists Hindi \emph{mol}. It survives in sayings and in \hi{मोलभाव}, bargaining.
\end{cousinweb}

\begin{grammarlens}[title={What you've built}]
One more word for getting around.
\end{grammarlens}

\subsection*{Your turn}
\begin{itemize}
\raggedright
  \item \emph{Say it:} \emph{mol}
  \item \emph{Say it:} \emph{mol}, once more
  \item \emph{Say it:} \emph{māṁ ke pyār kā koī mol nahīṃ hotā.} — A mother's love is beyond price.
  \item \emph{Recall:} say the Hindi for to collect (dues), then the Hindi for the going rate, then say \emph{mol} again
\end{itemize}

\begin{tcolorbox}[breakable,colback=teal!4,colframe=teal!35!black,title={Before you move on}]
What does \hi{मोल} mean? (\textquotedblleft{}Worth; a price (old-fashioned)\textquotedblright{}.) Say it once more.
\end{tcolorbox}
\section[\texorpdfstring{\hi{सौदा} (saudā)}{सौदा (saudā)}]{\hi{सौदा} (saudā) — a deal; shopping}

\label{lesson:HI-C431-sauda}



\begin{quote}
Before the new one: say the Hindi for the going rate, then the Hindi for worth.
\end{quote}

\subsection*{You'll want to know: \hi{सौदा}}
\textbf{\hi{सौदा}} — \emph{saudā} — \textquotedblleft{}a deal; shopping\textquotedblright{}.

\begin{quote}
\hi{हमने} \hi{दुकानदार} \hi{से} \hi{अच्छा} \hi{सौदा} \hi{किया।}
\end{quote}

\begin{quote}
\emph{hamne dukāndār se acchā saudā kiyā.}
\end{quote}

\begin{quote}
We got a good deal from the shopkeeper.
\end{quote}

Masculine.

\begin{cousinweb}
\textbf{\hi{सौदा}} (\emph{saudā}) is usually traced to Persian \emph{saudā}, \textquotedblleft{}trade, a bargain\textquotedblright{}; \hi{सौदा}-\hi{सुलफ़} is the day's shopping.
\end{cousinweb}

\begin{grammarlens}[title={What you've built}]
One more word for getting around.
\end{grammarlens}

\subsection*{Your turn}
\begin{itemize}
\raggedright
  \item \emph{Say it:} \emph{saudā}
  \item \emph{Say it:} \emph{saudā}, once more
  \item \emph{Say it:} \emph{hamne dukāndār se acchā saudā kiyā.} — We got a good deal from the shopkeeper.
  \item \emph{Recall:} say the Hindi for the going rate, then the Hindi for worth, then say \emph{saudā} again
\end{itemize}

\begin{tcolorbox}[breakable,colback=teal!4,colframe=teal!35!black,title={Before you move on}]
What does \hi{सौदा} mean? (\textquotedblleft{}A deal; shopping\textquotedblright{}.) Say it once more.
\end{tcolorbox}
\section[\texorpdfstring{\hi{सेल} (sel)}{सेल (sel)}]{\hi{सेल} (sel) — a sale (at reduced prices)}

\label{lesson:HI-C431-sel}



\begin{quote}
Before the new one: say the Hindi for worth, then the Hindi for a deal.
\end{quote}

\subsection*{You'll want to know: \hi{सेल}}
\textbf{\hi{सेल}} — \emph{sel} — \textquotedblleft{}a sale (at reduced prices)\textquotedblright{}.

\begin{quote}
\hi{कपड़ों} \hi{की} \hi{दुकान} \hi{में} \hi{सेल} \hi{लगी} \hi{है।}
\end{quote}

\begin{quote}
\emph{kapṛoṁ kī dukān meṁ sel lagī hai.}
\end{quote}

\begin{quote}
There is a sale on at the clothes shop.
\end{quote}

Feminine.

\begin{cousinweb}
\textbf{\hi{सेल}} (\emph{sel}) is borrowed from English \emph{sale}, a time of cut prices; it is feminine: \hi{सेल} \hi{लगी} \hi{है}.
\end{cousinweb}

\begin{grammarlens}[title={What you've built}]
One more word for getting around.
\end{grammarlens}

\subsection*{Your turn}
\begin{itemize}
\raggedright
  \item \emph{Say it:} \emph{sel}
  \item \emph{Say it:} \emph{sel}, once more
  \item \emph{Say it:} \emph{kapṛoṁ kī dukān meṁ sel lagī hai.} — There is a sale on at the clothes shop.
  \item \emph{Recall:} say the Hindi for worth, then the Hindi for a deal, then say \emph{sel} again
\end{itemize}

\begin{tcolorbox}[breakable,colback=teal!4,colframe=teal!35!black,title={Before you move on}]
What does \hi{सेल} mean? (\textquotedblleft{}A sale (at reduced prices)\textquotedblright{}.) Say it once more.
\end{tcolorbox}
\section[\texorpdfstring{\hi{किफ़ायती} (kifāyatī)}{किफ़ायती (kifāyatī)}]{\hi{किफ़ायती} (kifāyatī) — economical; good value}

\label{lesson:HI-C431-kifayati}



\begin{quote}
Before the new one: say the Hindi for a deal, then the Hindi for a sale.
\end{quote}

\subsection*{You'll want to know: \hi{किफ़ायती}}
\textbf{\hi{किफ़ायती}} — \emph{kifāyatī} — \textquotedblleft{}economical; good value\textquotedblright{}.

\begin{quote}
\hi{यह} \hi{होटल} \hi{किफ़ायती} \hi{है} \hi{और} \hi{साफ़} \hi{भी।}
\end{quote}

\begin{quote}
\emph{yah hoṭal kifāyatī hai aur sāf bhī.}
\end{quote}

\begin{quote}
This hotel is good value, and clean too.
\end{quote}

\begin{cousinweb}
\textbf{\hi{किफ़ायती}} (\emph{kifāyatī}) is built on \emph{kifāyat}, \textquotedblleft{}thrift\textquotedblright{}, usually traced to Arabic \emph{kifāya}, \textquotedblleft{}what is enough\textquotedblright{}, through Persian: spending no more than is enough.
\end{cousinweb}

\begin{grammarlens}[title={What you've built}]
One more word for getting around.
\end{grammarlens}

\subsection*{Your turn}
\begin{itemize}
\raggedright
  \item \emph{Say it:} \emph{kifāyatī}
  \item \emph{Say it:} \emph{kifāyatī}, once more
  \item \emph{Say it:} \emph{yah hoṭal kifāyatī hai aur sāf bhī.} — This hotel is good value, and clean too.
  \item \emph{Recall:} say the Hindi for a deal, then the Hindi for a sale, then say \emph{kifāyatī} again
\end{itemize}

\begin{tcolorbox}[breakable,colback=teal!4,colframe=teal!35!black,title={Before you move on}]
What does \hi{किफ़ायती} mean? (\textquotedblleft{}Economical; good value\textquotedblright{}.) Say it once more.
\end{tcolorbox}
